\documentclass[10pt,a4paper]{amsart}
\usepackage[margin=19mm]{geometry}
\usepackage{amsmath,amssymb,mathtools}
\usepackage[scr=rsfs,cal=euler]{mathalfa}
\usepackage{xfrac}
\usepackage[unicode,hidelinks,bookmarks=false]{hyperref}
\newcommand{\ie}{i.e.\ }
\newcommand{\cf}{cf.\ }
\newcommand{\resp}{respectively\ }
\newcommand{\Subsection}{\subsection}
\providecommand{\nobreakdash}{\nobreak\hbox{-}}
\setlength{\emergencystretch}{2em}
\multlinegap0pt
\usepackage{amssymb}
\usepackage{enumerate}
\usepackage{tikz}
\usepackage{dsfont}
\usepackage{xcolor}
\newtheorem{thm}{Theorem}
\newcommand{\e}{\varepsilon}
\title{Operator model and a trace formula for pairs of unitary operators}
\author{Daniel Alpay, Fabrizio Colombo, Izchak Lewkowicz, Irene Sabadini}
\date{Source: arXiv:2607.05334v1 (CC BY 4.0)}
\begin{document}
\maketitle

\noindent\emph{Source and licence.} Operator model and a trace formula for pairs of unitary operators. Version arXiv:2607.05334v1 (CC BY 4.0), distributed by arXiv under the Creative Commons Attribution 4.0 International licence, http://creativecommons.org/licenses/by/4.0/, as recorded in the arXiv licence metadata for this exact version and shown on the abstract page https://arxiv.org/abs/2607.05334v1. Full licence notice: \texttt{licenses/arxiv-2607.05334-CC-BY-4.0.txt}.\par\smallskip

\noindent\textit{Editorial scope.} Counted unit: the article's thm ththth (source0.tex lines 1144--1150). The article's complete original proof of it is delivered with it (source0.tex lines 1152--1247, one \texttt{proof} environment of 96 lines). No mathematics is written, repaired or re-derived: the statement and the proof are the article's own lines, in the article's own notation, and no retained line is substituted or re-flowed.  Retained uncounted as the source-local context the proof needs: lines 186--189; lines 215--218; lines 941--951; lines 1017--1020.  Nothing was omitted from any retained span, so the package carries no omission record.  No retained line contains a citation, so the wrapper carries no bibliography and no bibliography entry.  No retained line contains a figure environment, an image-inclusion command or drawing code, so no image is delivered and the package declares no asset dependency.  Editorial formatting only, in addition: an AMS \texttt{amsart} preamble that restates the article's own declarations and macros used by the retained lines, namely the environments \texttt{thm} on separate counters, together with the standard packages those lines need.  The retained environments print in this document as Theorem 1 in order of appearance, whereas the article prints the same items with the numbering of its own full document.  The complete native source of the article as distributed in the arXiv e-print for this exact version is bundled unchanged as \texttt{sources/204534/source0.tex}.
\begin{equation}
  \label{2-2}
\Phi({\lambda})=-\left(\Phi\left(\frac{1}{\overline{{\lambda}}}\right)\right)^*,\quad {\lambda}\in\mathbb E.
  \end{equation}

\begin{equation}
  \label{link}
\Phi(\lambda)=iE+\sum_{\ell=1}^mM_\ell \frac{e^{it_\ell}+\lambda}{e^{it_\ell}-\lambda},
\end{equation}

            \begin{thm}
              Let $\Phi$ be a $\mathbb C^{n\times n}$-valued rational function, analytic at infinity, with $\Phi(\infty)$ invertible (but $\Phi$ need not satisfy \eqref{2-2}).
              Let $\Phi({\lambda})=D+C({\lambda}I_N-A)^{-1}B$ be a realization of $\Phi$ centered at $\infty$. Then,
\begin{eqnarray}
  \label{eqas-1}
  \det\Phi({\lambda})&=&(\det D^{-1})\det\left(({\lambda}I_N-A^\times)({\lambda}I_N-A)^{-1}\right)\\
  {\rm Tr}\,\Phi^\prime({\lambda})\Phi({\lambda})^{-1}&=&{\rm Tr}\left\{({\lambda}I_N-A^\times)^{-1}-({\lambda}I_N-A)^{-1}\right\},
                                          \label{eqas-2}
\end{eqnarray}
where $A^\times=A-BD^{-1}C$.
            \end{thm}

\begin{equation}
\label{lala}
    {\rm Tr}\,( U_+^{*(\ell+1)}-U_-^{*(\ell+1)})={\rm Tr}\, CA^\ell B,\quad \ell=0,1,\ldots .
\end{equation}

        \begin{thm}
          \label{ththth}
          Assume $n=1$ and $\Phi$ rational and satisfies \eqref{2-2}. Then,
        \begin{equation}
          {\rm Tr}\,\,\left\{ f(U_-)-f(U_+)\right\}=\sum_{\lambda \,\,\mbox{\rm a pole or zero of $\Phi$}} f(\lambda).
        \end{equation}
      \end{thm}

      \begin{proof} We write
        \begin{equation}
          f(\lambda)=\sum_{\ell=0}^\infty f_\ell\lambda^\ell
        \end{equation}
        and divide the proof into a number of steps.\\

        STEP 1: {\sl It holds that
\[
  ( \lambda I_{\mathcal H}-U_-)^{-1}-({\lambda}I_{\mathcal H}-U_+)^{-1}=\sum_{\ell=0}^\infty
  \frac{1}{\lambda^{\ell+1}}\left(U_-^\ell-U_+^\ell\right),
\]
where the convergence is in norm for $|\lambda|>1$.\\
}

The next step is an improvement of \eqref{lala}, due to the meromorphicity of $\Phi$.\\

STEP 2: {\sl There exists a constant $c$ such that

\begin{equation}
  | {\rm Tr}\,(U_-^{\ell+1}-U_+^{\ell+1})|\le c ,\quad \ell=0,1,\ldots
  \label{lalala}
\end{equation}
}\\


Indeed, from \eqref{eqas-2} we have:\\

STEP 3: {\sl With  $R>1$ one has:
  \begin{equation}
    \frac{1}{2\pi i}\int_{|\lambda|=R}f(\lambda)\left(( \lambda I_{\mathcal H}-U_-)^{-1}-({\lambda}I_{\mathcal H}-U_+)^{-1}\right)d\lambda=f(U_-)-f(U_+).
    \label{trace-987}
  \end{equation}
}
Let $\e>0$. From \eqref{lalala} there exists $N$ such that
\begin{equation}
  \label{7-16}
\sum_{\ell=0}^\infty\frac{|{\rm Tr}\, (U_-^\ell-U_+^\ell)|}{R^{\ell+1}}<\e.
  \end{equation}
  We write
  \[
    \begin{split}
      \frac{1}{2\pi i}\int_{|\lambda|=R}f(\lambda)\left({\rm Tr}\,\left(( \lambda I_{\mathcal H}-U_-)^{-1}-({\lambda}I_{\mathcal H}-U_+)^{-1}\right)\right)d\lambda&=\\
      &\hspace{-6cm}=
        \frac{1}{2\pi i}\int_{|\lambda|=R}f(\lambda)\left({\rm Tr}\,
        \left(\sum_{\ell=0}^\infty \frac{1}{\lambda^{\ell+1}}\left(U_-^\ell
        -U_+^\ell\right)\right)\right)d\lambda\\                                                          &\hspace{-6cm}
                                                                                                            =\frac{1}{2\pi i}\int_{|\lambda|=R}\left({\rm Tr}\,\left(\sum_{\ell=0}^N
                                                                                                            \frac{f(\lambda)}{\lambda^{\ell+1}}\left(U_-^\ell-U_+^\ell
                                                                                                            \right)\right)\right.\\
      &\hspace{-3.5cm}\left.
        +\sum_{\ell=N+1}^\infty \left(\frac{f(\lambda)}{\lambda^{\ell+1}}\left(U_-^\ell
        -U_+^\ell\right)\right)\right)d\lambda\\
    &\hspace{-6cm}=\spadesuit+\clubsuit,\end{split}
\]
with
\[
  \begin{split}
    \spadesuit&=\frac{1}{2\pi i}\int_{|\lambda|=R}\left({\rm Tr}\,\left(\sum_{\ell=0}^N
                \frac{f(\lambda)}{\lambda^{\ell+1}}\left(U_-^\ell-U_+^\ell\right)
                \right)\right)d\lambda\\
              &=\frac{1}{2\pi i}\int_{|\lambda|=R}\left(\sum_{\ell=0}^N
                \frac{f(\lambda)}{\lambda^{\ell+1}}\left({\rm Tr}\,\left(U_-^\ell
                -U_+^\ell\right)\right)\right)d\lambda\\
              &=\sum_{\ell=0}^N\left(\frac{1}{2\pi i}\int_{|\lambda|=R}
                \frac{f(\lambda)}{\lambda^{\ell+1}}d\lambda\right)
                \left({\rm Tr}\,\left(U_-^\ell-U_+^\ell\right)\right)\\
              &={\rm Tr}\,\left(\sum_{\ell=0}^Nf_\ell U_-^\ell\right)-{\rm Tr}\,
                \left(\sum_{\ell=0}^Nf_\ell U_+^\ell\right).
              \end{split}
\]


Furthermore, by \eqref{7-16},
\[
  |\clubsuit|\le\e \max_{|\lambda|=R}|f(\lambda)|
\]
which ends the proof of the step since $\e$ is arbitrary.\\


STEP 4: {\sl To conclude the proof we compute
  \begin{equation}
    \label{residue-567}
\frac{1}{2\pi i}\int_{|\lambda|=R}f(\lambda)\frac{\Phi^\prime(\lambda)}{\Phi(\lambda)}d\lambda.
    \end{equation}
}

The poles of the function $f(\lambda)\frac{\Phi^\prime(\lambda)}{\Phi(\lambda)}$ are all simple and on the unit circle, as is seen from formula \eqref{link} and the similar
formula for $\Phi^{-1}$. They correspond to the eigenvalues of $\mathcal A^\times$
and $\mathcal A$. Thus each pole, say $w$, has residue



\[
          {\rm Res}_{z=w}f(z)\frac{\Phi^\prime(z)}{\Phi(z)}=f(w).
        \]
        \end{proof}

\end{document}
